\documentclass[12pt,english]{article}
\usepackage[T1]{fontenc}
\usepackage[latin9]{inputenc}
\usepackage[toc,page]{appendix}
\usepackage{geometry}
\geometry{verbose,tmargin=1.5in,bmargin=1.5in,lmargin=1in,rmargin=1in}
\usepackage{bm}
\usepackage{amstext}
\usepackage{comment}
\usepackage{amsmath}
\usepackage{hyperref}
\hypersetup{
     colorlinks = true,
     citecolor = navy,
     linkcolor = navy,
     urlcolor=navy
}
\usepackage[space]{grffile}
\usepackage{setspace}
\usepackage[normalem]{ulem}
\usepackage{xcolor}
\usepackage{soul}
\usepackage{dsfont}
\usepackage{refstyle}
\usepackage[capitalize,nameinlink]{cleveref}
\usepackage{setspace}
\usepackage{bbm}
\usepackage[tableposition=t]{caption} % ensures caption of longtables is normalsized.
\usepackage{mathpazo} % add Palatino font
\usepackage{longtable}
\usepackage{amssymb}
\usepackage{amsfonts}
\usepackage{amsmath}
%\usepackage[margin=1in]{geometry} % to set margins
\usepackage{scalefnt}
\usepackage{wrapfig}
\usepackage{graphicx}
%\usepackage{subfigure}
\usepackage{lscape}
\usepackage{pdflscape}
\usepackage{multirow}
\usepackage{array}
\usepackage{rotating}
\usepackage{bibunits}
\usepackage{fix-cm}
\usepackage{float}
\floatplacement{table}{H}
\usepackage{epsfig}
\usepackage{epstopdf}
\usepackage{refstyle}
\usepackage{booktabs}
\usepackage{tabularx}
\usepackage{array}

\newcolumntype{x}[1]{>{\raggedleft\arraybackslash}p{#1}}%
\newcolumntype{d}[1]{>{\centering\arraybackslash}p{#1}}%
\newcolumntype{y}[1]{>{\raggedright\arraybackslash}p{#1}}%
\usepackage{verbatim}
%\epstopdfsetup{suffix=}
\setcounter{MaxMatrixCols}{10}
\usepackage[longnamesfirst,authoryear]{natbib}
\newenvironment{proof}[1][Proof]{\noindent \textbf{#1.}}{\rule{0.5em}{0.5em}}
\setlength{\topmargin}{-0.80in} \setlength{\textheight}{9.50in}
\setlength{\evensidemargin}{.0in} \setlength{\oddsidemargin}{.0in}
\setlength{\textwidth}{6.85in}
\setlength{\voffset}{0.4cm}
\setlength{\hoffset}{-0.5cm}
\usepackage{setspace}
\usepackage{afterpage}
\usepackage[flushleft]{threeparttable}
\usepackage{tabulary}
%\usepackage[titletoc,title]{appendix}
\usepackage[bottom]{footmisc}
\usepackage{color}
	\definecolor{maroon}{RGB}{128,0,0}
	\definecolor{navy}{RGB}{0,0,165}
	\definecolor{lb}{RGB}{150,165,190}
	\definecolor{MyWhite}{RGB}{255, 255, 255}
	
\definecolor{red}{rgb}{0.9,0.1,0.1}
\definecolor{blue}{rgb}{0,0,0.7}
\definecolor{dgr}{rgb}{0.08,0.25,0.18}
\definecolor{navy}{RGB}{0,0,165}
\definecolor{org}{HTML}{EB811B}
\definecolor{lgr}{HTML}{14B03D}
\definecolor{shadecolor}{rgb}{0.8,0.8,0.8}

\definecolor{graphblue}{RGB}{64,130.56,194.5}
\definecolor{graphred}{RGB}{255,0,0}
\definecolor{graphyellow}{RGB}{237.8340,177.64,32}
\definecolor{graphpurple}{RGB}{126,47,142}

\newcommand{\hlight}[1]{\textcolor{org}{#1}}
\newcommand{\glight}[1]{\textcolor{lgr}{#1}}
\newcommand{\rlight}[1]{\textcolor{red}{#1}}
\newcommand{\blight}[1]{\textcolor{graphblue}{#1}}
\newcommand{\plight}[1]{\textcolor{graphpurple}{#1}}

\usepackage{colortbl}

\usepackage{caption}
\usepackage{subcaption}
\usepackage{placeins}
\usepackage{nicefrac}
%\usepackage{chngcntr}

\newcommand{\bit}{\begin{itemize}}
\newcommand{\eit}{\end{itemize}}
\usepackage{listings}
\def\cN{\mathcal{N}}
\newtheorem{theorem}{Theorem}
\newtheorem{remark}[theorem]{Remark}
\usepackage{mdframed}
\newmdtheoremenv{assumption}{Assumption}

\usepackage{enumitem}


\begin{document}


%\onehalfspacing


\title{Doing our supply/demand contribution graph properly}
\author{{\emph{Which influencer has the most influence?}}}
%Beggar or Prosper-thy-Neighbour? \\
%Fiscal Spillovers Across Firms and Countries
\maketitle

\paragraph{Purpose:} When I did the original note for the classification of sectors and being the \emph{originator} of a supply or demand shock, I thought we could do a similar graph to Shapiro's inflation graph showing contributions of demand and supply with our classifications directly. Based on our last meeting, it became clear that this was incorrect. If a sector (like oil) is the source of a supply shock \emph{everywhere} we need to have the right approach to capture this.

Here is what I would suggest in order to operationalize this. \hlight{Note that I will use the same notation of matrices that we have in our other notes -- $\mathbold{\Omega}$ is the share of intermediates row sources from column, $\underline{\gamma}$ is the intermediate share of sales for each industry.}

\begin{enumerate}
	\item Start by writing a function that produces the graphs. This function should take as inputs the following:
	\begin{itemize}
		\item $\pi_t$ -- CPI inflation ($T$ x $1$ total time series)
		\item $\mathbold{E}$ -- Expenditure weights by sector ($P$ x $T$) where $P$ is the number of products and $T$ is the number of time periods. Note that $1_P'E = 1_T'$ -- the sum of expenditure weights over the product/industry dimension $P$ is 1.
		\item A Demand shock matrix $\mathbold{D}$ ($P$ x $T$) where entries are between 0 and 1. If an entry is 0.75 for sector 3 in 2000Q1, then it means that in 2000Q1, sector 3 was 75\% affected by a demand shock and 25\% affected by a supply shock.
		\item A matrix of sectoral inflations $\mathbold{\Pi}$ ($P$ x $T$) where ($1_P'(\mathbold{E}\cdot \mathbold{\Pi})\approx \pi_t$) noting that $1_P$ is a $P$ x $1$ vector of ones and $A\cdot B$ means the elementwise multiplication of $A$ and $B$.\footnote{What I mean by this math is if you take the expenditure shares and scale them by inflation in each product/industry and then sum over products/industries, you should get CPI inflation and it should be close to the first input.}
	\end{itemize}
	Note that there might be more inputs I have forgotten. 
	
	This function will then plot CPI inflation as a line, then calculate $1_P'(\mathbold{D} \cdot \mathbold{E}\cdot \mathbold{\Pi})$ and plot it as a bar and call it the ``demand contribution to CPI''. For supply bars it will calculate $1_P'((1-\mathbold{D}) \cdot \mathbold{E}\cdot \mathbold{\Pi})$.
	
	\item Use this function to replicate the Shapiro classification. In this case entries of $D$ will either be 1 or 0.
	\item Next we will use this function to implement our classification. We will get a initial matrix $\hlight{\tilde{D}}$ from our procedure. This however only characterizes the sectors based on the \emph{initial} source of the shock (supply of demand). What we want is to calculate the influence of these shocks throughout the economy. To do this let's do the following:
	\begin{enumerate}
		\item We should already have a classification where if $P_{VA}$ (price of value added) and \underline{real} quantity (i.e. real production and not sales) both rise/fall in industry $i$ at time $t$, we say $\hlight{\tilde{D}_{i,t}}=1$. If instead $P_{VA}$ and quantity move in opposite directions, we say $\hlight{\tilde{D}_{i,t}}=0$.\footnote{Tony can you please check we are doing the classification this way?}
		\item Next we want to calculate the influence of this onto prices.\footnote{Ben, note that if we wanted to do the influence on quantities we would need to do something a little different to this. Let me know if you want the math for that.} To find that first let us calculate the following $\mathbold{\Theta}\equiv (I-diag(\underline{\gamma})\mathbold{\Omega})^{-1} = (I-\mathbold{\Delta})^{-1}$. \rlight{Note that $\mathbold{\Theta}\neq \mathbold{\Xi}$ because $\mathbold{\Xi}=(I-\mathbold{\Delta}')^{-1}$}. $\mathbold{\Theta}$ is also a $P$ x $P$ matrix.
		\item Next we should calculate the influence shares of $P_{VA}$ movements onto output prices (and therefore sectoral inflation). However doing so is a little tricky because even though $\mathbold{\Theta}$ is constant, the value added prices vary over time. 
		\item Next calculate $\mathbold{\theta}_t\equiv \mathbold{\Theta}\cdot (\log(\underline{P}_{VA,t})\otimes 1_P')$ where $\underline{P}_{VA,t}$ is a $P$ x $1$ matrix when doing this multiplication.\footnote{$\otimes $ represents the kronecker product but effectively all I want is $P$ repetitions of the vector $\log(\underline{P}_{VA,t})$ so that all rows have the same number in them.} Note that we need to do this for every time period $t=[1,\dots,T]$. Note also that $\mathbold{\theta}_t$ is a $P$ x $P$ matrix (and the we have $T$ of these).
		\item Calculate the column sum of $\underline{\kappa}_t = 1_P'\mathbold{\theta}_t$. Note that $\underline{\kappa}_t$ should be $1$ by $P$.
		\item Influence is calculated by $\mathbold{\lambda}_t\equiv diag(1/\underline{\kappa}_t)\mathbold{\theta}_t$. This is a $P$ x $P$ matrix and there are $T$ of them -- one for each $t$. Since $\mathbold{\Lambda}_t$ is an influence share, the column-sum should be 1: $1_P'\mathbold{\lambda}=1_P$ or summing over rows within a column should be 1.
		\item Calculate the input $\mathbold{D}$ one column at a time. That is if $\underline{D}_t$ is the $t^{th}$ column of $\mathbold{D}$ and $\hlight{\underline{\tilde{D}}_t}$ is the $t^{th}$ column of $\hlight{\tilde{D}}$, then it is calculated as: $\underline{D}_t\equiv \mathbold{\lambda}_t\hlight{\tilde{\underline{D}}_t}$. Note that since $\mathbold{\lambda}_t$ is $P$ x $P$ and $\hlight{\tilde{D}_t}$ is $P$ x $1$, each $\underline{D}_t$ is $P$ x $1$. 
		\item Finally concatenate the $\underline{D}_t$s to get:
		
		$$
			\mathbold{D} = \begin{bmatrix}
				\underline{D}_1 &\dots &\underline{D}_T
			\end{bmatrix}
		$$
	\end{enumerate}
	\item Before using $\mathbold{D}$ as an input we should verify that all entries are between 0 and 1. Let me know if any are negative -- this might be possible and we need to think about what to do in that case.
\end{enumerate}


\end{document}